%% kapa: exam=1 round=20 year=2009 period=1 form=A subject=civil_law subject_ko=민법 seq=38 number=38
%% source: Q-Net 감정평가사 기출문제 (cst00309) artlSeq=5024178 file=extracted/1008/5024178/2009년 제20회 감정평가사 1차시험 기출문제/1교시.hwp
%% conversion: kordoc 4.24.51 -> tools/gen1.py (auto); review=proofread-sonnet; flags=none
\begin{question}{38}
甲이 채무자 乙소유의 A부동산에 채권최고액 2천만원의 근저당권을 취득하였다. 다음 중 옳지 \uline{않은} 것은? (다툼이 있으면 판례에 의함)
\begin{choices}
  \item 甲이 경매를 신청하는 경우 경매신청시에 근저당권이 확정되므로 그 이후에 발생하는 甲의 원금채권은 근저당으로 담보되지 않는다.
  \item 甲이 경매를 신청하여 경매개시결정이 있은 후에 경매신청이 취하되었다면 채무확정의 효과는 번복된다.
  \item 근저당권의 실행비용은 甲의 채권최고액 2천만원에 포함되지 아니한다.
  \item 甲의 채권이 확정되기 전이라도 피담보채권과 함께 근저당권을 양도할 수 있다.
  \item 만일 丙이 A부동산에 대하여 후순위 저당권을 취득한 경우, 丙이 경매를 신청하는 경우에는 매각대금 완납시에 甲의 근저당권이 확정된다.
\end{choices}
\end{question}
\answer{2}
